\documentclass[a4paper,11pt]{article}
\usepackage[margin=1in]{geometry}
\usepackage{amsmath, amssymb, amsthm}
\usepackage{graphicx}
\usepackage{subcaption}
\usepackage{booktabs}
\usepackage{multirow}
\usepackage{float}
\usepackage[table]{xcolor}
\usepackage{enumitem}
\usepackage{wrapfig}
\usepackage{listings}
\usepackage{hyperref}   % load last
\bibliographystyle{plain} 
\newtheorem{theorem}{Theorem}[section]
\newtheorem{proposition}{Proposition}[section]
\newtheorem{lemma}[theorem]{Lemma}

\hypersetup{colorlinks=true, linkcolor=red, citecolor=red, urlcolor=magenta}

\lstset{
  basicstyle=\ttfamily\small,
  breaklines=true,
  frame=single,
  showstringspaces=false,
  keywordstyle=\color{blue}\bfseries,
  commentstyle=\color{teal!70!black},
  stringstyle=\color{red!70!black}
}
\title{\textbf{Laplace Analysis and Applications}}
\author{%
  Pierre-Simon Laplace\textsuperscript{1}, Oliver Heaviside\textsuperscript{2}, and Gustav Robert Kirchhoff\textsuperscript{3} \\
  \small
  \textsuperscript{1}Acad´emie des Sciences, Paris, France\\
\textsuperscript{2}Royal Society, London, United Kingdom\\
\textsuperscript{3}University of Heidelberg, Heidelberg, Germany
}
\date{}

% \author{Johannes Kepler\thanks{Imperial Mathematician, Prague}}
\begin{document}



\maketitle
\section{From Functions to Transformations}
\subsection{The Laplace Transform}
Let $f(t)$ be a function defined for $t\,\geq 0$. Its Laplace transform is defined by
\begin{equation}
    \mathcal{L}\{f(t)\}=F(s) \int_0^\infty e^{-st}f(t)\, dt, s\gt 0
\end{equation}
The transform changes a function of time into a function of the new variable s.\footnote{The variable $s$ is commonly interpreted as a complex-frequency variable.}Symbolically,
\begin{equation*}
    f(t) \xrightarrow{\mathcal{L}} F(s).
\end{equation*}
The inverse Laplace transform reverses this process:
\begin{equation}
    f(t)=\mathcal{L}^{-1} \{ F(s)\}.
\end{equation}
The main purpose of the method is to replace differentiation and integration in the time domain
by simpler algebraic operations in the transform domain \cite{kreyszig2011}. A concise overview is available from
\href{https://mathworld.wolfram.com/LaplaceTransform.html}{the Laplace Transform entry on MathWorld}
\subsection{A Simple Example}
For the constant function $f(t)=1$,
\begin{align}
    \mathcal{L}\{ 1\} &=\int_0^\infty e^{-st}\ \notag
\\    &=\left[ -\frac{1}{s}e^{-st}\right]_0^\infty \notag \\
&=\frac{1}{s}
\end{align}
\newpage
\section{Common Transform Equations}
Table \ref{tab:laplace} lists several frequently used Laplace-transform equations.
\begin{table}[H]
\centering
\begin{tabular}{|c|c|c|c|}
\hline
\multirow{2}{*}{\textbf{Family}} &
  \multicolumn{2}{c|}{\textbf{Transform Equation}} &
  \multirow{2}{*}{\textbf{Condition}} \\
\cline{2-3}
 & \textbf{Time function} & \textbf{Laplace Transform} & \\
\hline
\multirow{3}{*}{Algebraic}
  & 1 & $\frac{1}{s}$ & $$s\gt0$$ \\
 & t & $\frac{1}{s^2}$ & $$s\gt0$$ \\
 & $t^n$ & $\frac{n!}{s^{n+1}}$& $n=0,1,2...$ \\
\hline
\multirow{2}{*}{Exponential}
  & $e^{\alpha t}$ & $\frac{1}{s-a}$ &   \\
 & $e^{-\alpha t} & $\frac{1}{s+a}$ & $$s \gt -\alpha$$ \\
\hline
\multicolumn{2}{|l|}{\textbf{Special Row Label}}
  & special-value & special-cond \\
\hline
\end{tabular}
\caption{The formulas in Table 1 allow many functions to be transformed without repeatedly evaluating
improper integrals.}
\label{tab:laplace}
\end{table}


\section{Application to an RC Circuit}
Figure \ref{fig:rc} shows a simple resistor--capacitor
circuit connected to an input voltage source.
\begin{wrapfigure}{r}{0.36\linewidth}
    \centering
    \includegraphics[width=\linewidth]{rc_circuit.png}
    \caption{Wrapped image.}
    \label{fig:rc}
\end{wrapfigure}
The resistor limits the current, while the capacitor stores electrical energy.
For a resistor R and capacitor C, Kirchhoff’s
voltage law gives
\begin{equation}
    RC\frac{dv_C(t)}{dt}+v_C(t)=v_{\text{in}}(t),
\end{equation}
where $v_C(t)$ is the voltage across the capacitor.
Assuming $v_C(t)=0$, taking the Laplace
transform of Equation \eqref{eqn:5} gives
\begin{align}
    &RC[sV_C(s)-v_C(0)]+V_C(s) \;\;=V_{\text{in}}(s)\notag \\
    &\implies RCsV_C(s)+V_C(s)\;\;=V_{\text{in}}(s) \notag \\
    &\implies V_C(s)(RCs+1)\;\;=V_{\text{in}}(s)
\end{align}
The transfer function is therefore
\begin{equation}
    H(s)=\frac{V_C(s)}{V_\text{in}(s)}=\frac{1}{RCs+1}
\end{equation}
For a constant input $V_0$
\begin{equation*}
    {V_\text{in}(s)}=\frac{V_0}{s}
\end{equation*}
Thus,
\begin{equation}
    V_C(s)=\frac{V_0}{s(RCs+1)}
\end{equation}
Taking the inverse transform gives
\begin{equation}
   v_C(t)=V_0(1-e^{-\frac{t}{RC}})
    \label{eq:9}
\end{equation}
The quantity $\tau=RC$ is called the time constant of the circuit. A larger value of $\tau$ causes the
capacitor to charge more slowly.\cite{ogata2010}
\section{How Laplace Analysis Is Used}
The Laplace transform is useful in many areas:
Differential equations  initial conditions are included automatically;
 derivatives become algebraic expressions.
Electrical circuits  voltages and currents are studied in the s-domain;
components lead to transfer functions.
Control and signal systems  stability can be studied through poles;
 common test inputs include:
the unit-step input;
the impulse input;
sinusoidal inputs.

\textcolor{blue}{In the time domain, a system may be described by a differential equation}, whereas \textcolor{red}{in the Laplace
domain, the same system is often described by a rational algebraic expression}.
\section{A Short Python Implementation}
Listing 1 evaluates the capacitor voltage at a finite set of time instants using an external Python
file.
\newpage
\lstinputlisting[language=Python,
  caption={ Computing samples of the RC step response.},
  label={lst:external}]{rc_response.py}
Equation \eqref{eq:9} and Listing \ref{lst:external} describe the same computation.
\bibliography{laplace}
\end{document}
